\documentclass[oneside,final,14pt,a4paper]{extreport}


%=====================================================================
%  BACHELOR THESIS PROPOSAL (REVISED after supervisor feedback)
%=====================================================================


\usepackage[T1]{fontenc}
\usepackage[utf8]{inputenc}
\usepackage[english]{babel}
\usepackage{tempora}

% Monospaced font: Courier (bold available via psnfss).
% The former manual \DeclareFontShape{OT1}{cmtt}{bx}{n}{...} was dead code:
% with \ttdefault=pcr it was never used, so it has been removed.
\renewcommand{\ttdefault}{pcr}


\usepackage{vmargin}
\setpapersize{A4}
\setmarginsrb{2.5cm}{2.2cm}{2.2cm}{2.2cm}{0pt}{10mm}{0pt}{13mm}
\usepackage{setspace}
\setstretch{1.5}
\usepackage{indentfirst}
\parindent=1.25cm


\usepackage{amsmath,amsfonts}
\usepackage{amsthm}
\usepackage{upgreek}
\usepackage{bm}


\usepackage{xcolor}
\usepackage{graphicx}
\graphicspath{{figs/}}
\usepackage{float}
\usepackage{array}
\usepackage{multirow}
\usepackage{booktabs}
\usepackage{longtable}
\usepackage[font=singlespacing,labelfont=bf]{caption}
\usepackage{subcaption}
\usepackage{enumitem}


\usepackage{tikz}
\usetikzlibrary{positioning,arrows.meta}


\usepackage{csquotes}
\usepackage[backend=biber,style=ieee,autocite=inline]{biblatex}
\addbibresource{references.bib}
\DefineBibliographyStrings{english}{%
  bibliography = {References},}


\usepackage{chngcntr}
\usepackage{pdfpages}
\usepackage{placeins}
\usepackage{fancyhdr}


\captionsetup[table]{labelfont={normalfont}, name={TABLE}, labelsep={newline}}
\captionsetup[figure]{labelfont={normalfont}, name={Fig.}, labelsep=period}


\renewcommand{\thesection}{\arabic{section}}
\counterwithin{figure}{section}
\counterwithin{table}{section}
\counterwithin{equation}{section}
\setcounter{secnumdepth}{3}


\setlength{\headheight}{32.0976pt}

% Modern fancyhdr interface (replaces the deprecated fancyplain style,
% preserving the original header layout for this oneside document:
% left header = section mark, right header = page number, empty footer).
% \pagestyle{fancy}
% \fancyhf{}
% \renewcommand{\sectionmark}[1]{\markright{\thesection\ #1}}
% \fancyhead[L]{\bfseries\rightmark}
% \fancyhead[R]{\bfseries\thepage}
% \renewcommand{\headrulewidth}{0.4pt}


\usepackage{hyperref}
\hypersetup{colorlinks=true, linkcolor=black, citecolor=black, urlcolor=black}


\begin{document}

\begingroup
\setlength{\parindent}{0pt}
\setlength{\parskip}{0pt}

\begin{center}
    {\Large\bfseries
    Confidence-Aware Adaptive Verification of LLM Recommendations\\
    in Introductory Programming Tasks\par}
    \vspace{0.5em}
    {\normalsize Student: Lidia Davydova\par}
    \vspace{0.15em}
    {\normalsize Supervisor: Paolo Ciancarini\par}
\end{center}

\vspace{0.5em}

\begin{center}
\begin{minipage}{0.85\textwidth}
\small
\noindent\textbf{Abstract.}
Large language models (LLMs) are increasingly used in programming
education, but their recommendations can be plausible yet incorrect,
so the quality of human--AI collaboration depends on
\emph{appropriate reliance}: accepting correct recommendations and
rejecting incorrect ones. Verification support can help users evaluate
recommendations, and introductory programming is well suited to it
because correctness can be checked through independent, execution-based
evidence such as tests, traces, and counterexamples. Providing such
evidence on every trial, however, introduces unnecessary interaction
cost. This thesis investigates whether a simple, interpretable adaptive
policy---gating verification availability on the participant's initial
confidence---can reduce verification burden relative to always-available
verification while preserving decision quality. A within-subject
experiment (provisional $N=60$) compares no-verification, static, and
adaptive conditions on a benchmark of 18 introductory Python tasks with
frozen, pre-validated Qwen3-8B recommendations balanced between correct
and incorrect. Trial-level outcomes are analysed with mixed-effects
logistic regression with participant and task random effects;
verification burden is measured through availability, examination,
and time measures and reported separately from quality outcomes.
Expected contributions include trial-level evidence on the
burden--quality trade-off of confidence-gated verification, a validated
reproducible benchmark with an instrumented prototype, and a complete
analysis pipeline.
\end{minipage}
\end{center}

\vspace{0.5em}
\endgroup

%=====================================================================
\section{Introduction and Research Problem}
%=====================================================================


Large language models (LLMs) are increasingly used in programming education
and software engineering to generate, explain, and debug code and to propose
answers to tasks \cite{Ma2024,Cambaz2024,Denny2024}. Because such systems can
produce plausible but incorrect outputs, the quality of the joint human--AI
outcome depends not only on model accuracy but on \emph{appropriate reliance}
\cite{LeeSee2004}: accepting correct recommendations and rejecting incorrect
ones. Over-reliance and under-reliance are the corresponding failures
\cite{PoursabziSangdeh2021,Vasconcelos2022}.


Verification support can help users evaluate recommendations, and programming
tasks suit it well: correctness can often be checked independently through
execution, tests, traces, or counterexamples---evidence grounded in the task
itself rather than in another LLM-generated justification. This matters
because explanations assist only when users can interpret them and connect
them to the recommendation, and simply providing more information does not
automatically reduce over-reliance
\cite{Bucinca2021,Vasconcelos2022,Bo2025,Schemmer2022}. Moreover, one can
only verify an answer one has the grounding to check \cite{Zhang2024Verify};
introductory programming provides exactly such grounding. However,
verification on every trial introduces unnecessary interaction cost, which
motivates providing verification evidence \emph{selectively}
\cite{CascadedLM2025}.


\subsection{Research Problem}

The research problem is how to provide verification support selectively, using
only information available before the correctness of the LLM recommendation is
known. The thesis investigates whether a simple adaptive policy, based
primarily on the user's initial confidence, can provide independent
verification support at lower burden than always-available verification, and
how verification burden trades off against decision quality. The focus is a
simple, interpretable policy---not a new LLM, a new verification algorithm, or
an optimisation method. The contribution is therefore not the threshold rule
itself; it is the controlled empirical evaluation of whether a simple
confidence gate can reduce verification exposure while preserving appropriate
reliance compared with always-available and no-verification baselines.


%=====================================================================
\section{Research Questions}
%=====================================================================


\textbf{Primary RQ.} How does confidence-aware adaptive verification affect
reliance on LLM recommendations in introductory programming tasks compared
with always-available (static) verification and no verification?

\textbf{RQ1.} How do no, static, and adaptive verification affect final
answer accuracy, over-reliance, and under-reliance?

\textbf{RQ2.} What trade-off does adaptive verification exhibit between
verification burden (verification availability, evidence-examination rate,
evidence-examination time) and outcome quality (final accuracy, appropriate
reliance) relative to static verification?

\textbf{RQ3.} How do the burden and effectiveness of the adaptive policy
differ across initial confidence levels and task types?

Neither task type nor task difficulty is an input to the adaptive policy
(Section~\ref{subsec:policy}); both are analysed only as exploratory
moderators of policy effectiveness.

A policy premise check examines whether initial confidence is associated with
initial answer correctness (Section~\ref{subsec:premise}).


%=====================================================================
\section{Novelty and Contributions}
%=====================================================================


The novelty is \emph{not} the observation that confidence relates to reliance
on AI---miscalibration is well documented \cite{KrugerDunning1999,Bo2025} and
taken here as a premise---but the combination of elements usually studied in
isolation: (i)~confidence as a per-trial \emph{trigger} deciding \emph{when}
support is shown, rather than an outcome to be calibrated or an AI-side
uncertainty signal; (ii)~\emph{independent, execution-based} evidence (tests,
traces, counterexamples) rather than additional AI-generated explanations;
(iii)~a programming setting in which recommendations are pre-validated and
objective evidence exists; and (iv)~an explicit burden--quality trade-off
evaluation, treating verification burden as a first-class outcome
\cite{CascadedLM2025,Zhang2024Verify}.


The thesis is expected to contribute: (1)~trial-level empirical evidence,
analysed with mixed-effects models, on how a simple confidence-based policy
trades verification burden against accuracy and appropriate reliance
(RQ1--RQ3); (2)~a validated benchmark of 18 introductory Python tasks with
fixed, pre-validated LLM recommendations (balanced correct/incorrect),
task-type and difficulty labels, and verification data; (3)~a reproducible
prototype of the three conditions with complete logging, including evidence
examination; (4)~an analysis pipeline for all outcome and burden measures.
The contribution is deliberately not a new LLM, a new algorithm, or a
tutoring system: the confidence gate is intentionally simple, and its value
lies in the controlled empirical comparison against always-available and
no-verification baselines in a reproducible human--AI interaction experiment.

%=====================================================================
\section{State-of-the-Art}
%=====================================================================

\paragraph{Trust and appropriate reliance.}
Trust in automation is classically defined as the attitude that an agent will
help an individual achieve a goal under uncertainty and in risky situations,
and the design goal is \emph{appropriate} reliance rather than maximal trust
\cite{LeeSee2004}. In AI-assisted decision-making, users frequently fail to
calibrate reliance to recommendation quality: over-reliance on incorrect
advice and under-reliance on correct advice persist even when additional
information about the model is available
\cite{PoursabziSangdeh2021,Chen2023}. Comparable miscalibration between
confidence and correctness is documented for human judgment itself
\cite{KrugerDunning1999}.

\paragraph{Explanations and cognitive forcing.}
Cognitive forcing functions, which compel deliberate engagement with AI
advice before it is accepted, can reduce over-reliance \cite{Bucinca2021}.
However, meta-analytic evidence finds no reliable benefit of explanations
over AI predictions alone \cite{Schemmer2022}, and experimental work shows
that explanations reduce over-reliance only when they are interpretable and
help users anticipate AI errors \cite{Vasconcelos2022}; interventions
evaluated with LLM advice similarly reduced over-reliance without improving
appropriate reliance overall \cite{Bo2025}. Because explanation-based
support may itself provide additional AI-generated information, this
motivates verification support grounded in independent, execution-based
evidence rather than additional AI text.

\paragraph{Verification and its cost.}
Correctness alone is an insufficient evaluation criterion: human scrutiny of
AI outputs is valuable but consumes time and effort, which motivates
frameworks that invoke human input selectively \cite{CascadedLM2025}.
Furthermore, one can only verify an answer one has the grounding to check
\cite{Zhang2024Verify}; introductory programming supplies such
grounding---executable reference answers, tests, and traces---so objective
verification evidence is available at decision time.

\paragraph{Confidence as a signal for adaptive support.}
Closest to this thesis, \cite{Bo2025} evaluate interventions for reliance on
LLM advice and find that interventions reduce over-reliance but generally
fail to improve appropriate reliance, while users' confidence is poorly
calibrated. The authors discuss confidence calibration as a design
implication. Related work adapts support to the user's trust state
\cite{SrinivasanAdjustForTrust}. Here, confidence is instead an
\emph{input signal}: a lightweight trigger deciding \emph{when} independent
support is shown, leaving the user's internal state unchanged. The documented
relationship between confidence and correctness is taken as a premise: if
confidence carries information about correctness, it can also inform when
verification is worth its interaction cost---a burden dimension absent from
calibration-oriented work and a primary outcome here (RQ2).

\paragraph{LLM assistance in programming and research gap.}
LLMs are increasingly used for code generation, debugging, and explanation in
computing education \cite{Ma2024,Cambaz2024,Denny2024}, but can be plausible
yet incorrect, so users must evaluate recommendations critically. Prior work
establishes that intuition, explanation design, and decision difficulty shape
reliance \cite{Chen2023,Vasconcelos2022,Zhang2024Verify}; what remains
unknown is whether a participant's initial confidence can serve as a
lightweight trigger for selectively displaying \emph{independent verification
evidence} after an LLM recommendation in introductory programming
tasks---the gap addressed here.


%=====================================================================
\section{Goal of the Thesis}
%=====================================================================

The goal of the thesis is to design, implement, and empirically evaluate an
experimental system that provides selective verification support for LLM
recommendations in introductory programming tasks. The study investigates the
trade-off between decision quality and observable verification burden produced
by a simple confidence-based adaptive policy relative to an always-available
(static) policy.

The project schedule is provided in Section~\ref{sec:workplan}.

\subsection{Task Domain and Benchmark}
\label{subsec:answerspace}

The benchmark comprises 18 introductory Python programming tasks 
of three types: \emph{identify bug}, \emph{check against specification}, 
and \emph{select correct implementation}. ``Write or repair code'' 
tasks are excluded because they do not provide a sufficiently constrained 
response space for the primary reliance measures.

The benchmark metadata include task type, difficulty level, reference answer,
fixed LLM recommendation, recommendation correctness, and task-specific
verification data. The 18 LLM recommendations are balanced between nine
correct and nine incorrect recommendations. Correct and incorrect
recommendations are distributed across the three task types as evenly as the
benchmark construction permits.

\emph{Task difficulty.} Each task is assigned a difficulty level
(\emph{low}, \emph{medium}, or \emph{high}) before the main experiment
according to a predefined rubric based on code complexity, the number of
reasoning steps required, and the programming constructs involved. Difficulty
ratings are checked during the pilot against pilot accuracy and response
times and are frozen together with the benchmark. Difficulty is not used as
an input to the adaptive policy and is analysed only as an exploratory
factor; if the pilot indicates that the ratings are unreliable, difficulty is
dropped from the confirmatory analysis and retained only as descriptive
metadata.

Each task has a closed-form response space that makes correctness 
and changes in the participant's answer well-defined. For \emph{identify bug}, 
the response consists of a binary bug-presence judgment and, when a bug 
is present, a line number identifying its location. Optional free-text 
explanations are recorded only for qualitative analysis. 
For \emph{check against specification}, the response is a binary 
\texttt{COMPLIANT}/\texttt{NOT\_COMPLIANT} judgment. 
For \emph{select correct implementation}, the response is an identifier 
corresponding to one of a fixed set of candidate implementations.

Each participant completes 18 trials, with six trials assigned to each
experimental condition. Each condition contains two tasks of each of the
three task types. Within each condition, three trials contain correct LLM
recommendations and three trials contain incorrect LLM recommendations, with
one correct and one incorrect recommendation for each task type.

Because condition membership varies within task across participants,
task-level heterogeneity is not systematically confounded with the
experimental condition, allowing task to be included as a random effect in
the analysis models.

\subsection{LLM Recommendation Specification}
\label{subsec:llm-spec}

All LLM recommendations are generated offline before the main experiment. No
LLM generation is performed during participant trials. Consequently, all
participants receive the same frozen recommendation for a given task, and
differences between experimental conditions cannot be attributed to runtime
LLM sampling variance.

\emph{Model and version.}

The recommendation generator is a fixed local snapshot of a specific
\texttt{Qwen3-8B} model artifact, an 8.2B-parameter instruction-tuned model
distributed by Qwen. The experiment uses an 8-bit quantized model in GGUF
format (\texttt{Q8\_0}). The exact model file, model revision, quantization,
inference backend, software version, and model-file SHA-256 hash are recorded
and frozen before benchmark generation. The benchmark therefore refers to a
specific frozen model artifact rather than to the mutable upstream
\texttt{Qwen3-8B} model name.

\emph{Inference configuration.}

Qwen3 is used in non-thinking mode so that generated recommendations contain
only the requested answer and do not depend on a model-generated reasoning
trace. The generation configuration is fixed as follows:

\begin{verbatim}
model: Qwen3-8B model artifact, 8-bit Q8_0
thinking: disabled
temperature: 0.7
top_p: 0.8
top_k: 20
max_tokens: 256
frequency_penalty: 0
\end{verbatim}

The sampling parameters \texttt{top\_k} and \texttt{frequency\_penalty} are
backend-dependent; their concrete semantics are those of the archived backend
version.

\emph{Reproducibility.}

The exact inference backend and version, command-line options, context
length, and hardware configuration are archived with the benchmark. The
generation seed is archived when supported by the inference backend. Because
GPU inference may be nondeterministic even under a fixed seed, reproducibility
of the benchmark is defined with respect to the archived model artifact and
the archived generation outputs rather than bit-exact regeneration.

\emph{Hardware configuration}
\begin{verbatim}
VRAM: Tesla V100 16 GB
RAM: 128 GB
\end{verbatim}

\emph{Prompt template.}

For each task, the model receives a structured prompt containing:
(i) the task type and instructions, (ii) the task description and any Python
code snippet, (iii) a strict output schema, and (iv) an instruction to return
only the requested answer.

The prompt does not contain information about the experimental condition,
verification policy, participant confidence, or intended correctness of the
recommendation. The complete prompt template and all task-specific prompts are
archived with the benchmark.

\emph{Candidate generation.}

For each benchmark task, multiple candidate LLM recommendations are 
generated using a fixed inference configuration. The model receives 
a task-specific prompt containing the task description, required instructions, 
and a strict output format corresponding to the task's predefined response 
space. The model is instructed to solve the task and return only the 
requested answer, without explanations or confidence estimates.

The LLM recommendation is a task-specific structured answer expressed 
in the predefined response space of the task. For example, 
for an \emph{identify bug} task, it 
may be a predefined bug location or type (e.g., \texttt{LINE\_3}); 
for a \emph{check against specification} task, it may be a binary judgment 
such as \texttt{SATISFIES} or \texttt{DOES\_NOT\_SATISFY}; and for a 
\emph{select correct implementation} task, it may be an implementation 
identifier such as \texttt{B}. The exact response representation is defined 
separately for each task before candidate generation.

Each model output is parsed into the task's canonical structured response 
representation. Outputs that cannot be parsed unambiguously are rejected 
and regenerated. Candidate recommendations are then evaluated against the 
reference answer to determine their correctness.

For each task, one recommendation is selected for the final benchmark. 
The benchmark contains an equal number of correct and incorrect recommendations. 
Correct recommendations must correspond to the reference answer. Incorrect 
recommendations must be objectively incorrect while remaining sufficiently 
plausible to a novice participant. Selection is performed before the participant 
study and does not depend on participant data or experimental condition.

The final recommendation is fixed for the entire study: the same task always 
uses the same LLM recommendation regardless of experimental condition or 
participant. The recommendation is stored in canonical structured form for 
analysis and rendered in a participant-facing form appropriate to the task type.

For correct-recommendation tasks, one valid candidate is selected. For
incorrect-recommendation tasks, an incorrect candidate is selected only if it
is both objectively incorrect and sufficiently plausible to a novice
participant. If no suitable incorrect candidate is produced by the model, the
benchmark construction protocol permits a documented minimal perturbation of
a correct answer, such as an off-by-one value, a wrong operator, or another
task-appropriate error. Such perturbations are recorded separately from
model-generated candidates.

\emph{Validation procedure.}

Recommendation correctness is first established automatically whenever an 
objective reference procedure is available. Depending on the task type, this 
may involve program execution, reference-output comparison, or predefined tests. 
When automatic validation is not applicable, correctness is established by comparison 
with the predefined reference answer.

The resulting correctness label and the plausibility of incorrect 
recommendations are then independently reviewed by two annotators with 
introductory programming expertise. Annotators do not receive information 
about the experimental condition or the adaptive policy. Disagreements are 
resolved by a third annotator according to the predefined validation procedure.

An incorrect recommendation is retained only when it is objectively incorrect 
and judged sufficiently plausible that its incorrectness is not immediately 
obvious to a novice participant. The validation procedure, candidate pool, 
selected candidate, correctness label, annotator decisions, and any 
perturbation are archived.

After validation, the benchmark is frozen. No recommendation, correctness 
label, prompt, model parameter, or candidate-selection decision is changed 
during the main experiment. Recommendation correctness is hidden from 
participants and is never provided to the adaptive policy. The correctness
label is stored only in the analysis dataset and is not available to the
participant-facing application or to the policy module during a trial.


\subsection{Verification Support and Burden Instrumentation}
\label{subsec:burden}

For each task, the verification mechanism is fixed before the experiment 
and is identical across experimental conditions. The only manipulated factor 
is whether the verification evidence is made available to the participant.

Verification is task-specific and provides independent evidence that 
can be used to evaluate the LLM recommendation. For \emph{identify bug} tasks, 
the system provides predefined tests with failing cases or concrete 
counterexamples that provide evidence about the presence and location of a bug. 
For \emph{check against specification} tasks, the system provides 
specification-based test cases with expected and actual results. For 
\emph{select correct implementation} tasks, the candidate implementations are 
executed against a predefined test suite, and the resulting test outcomes are 
displayed.

Verification artifacts are constructed during benchmark development 
by the researchers and are derived independently from the reference answer 
and task specification. They are not generated by the LLM used to produce 
the experimental recommendations. Each verification artifact has a predefined 
expected outcome and is validated before the main experiment.

For \emph{identify bug} tasks, the researchers construct a set of targeted 
test cases or counterexamples designed to expose the specified bug and provide 
evidence about its location or effect. For \emph{check against specification} 
tasks, the researchers construct specification-based test cases covering the 
relevant requirements, with expected results derived from the task 
specification. For \emph{select correct implementation} tasks, the 
researchers construct a predefined test suite containing cases that 
distinguish the candidate implementations, with expected outcomes derived 
from the task requirements.

Each verification artifact is evaluated against the reference answer or 
reference implementation during benchmark validation. The verification 
artifacts are then independently reviewed by two annotators with introductory 
programming expertise to confirm that they are correct, relevant to the task, 
and sufficient to provide evidence for the intended judgment. Disagreements 
are resolved by a third annotator according to the predefined validation 
procedure.

After validation, the verification artifacts are frozen together with the 
benchmark. No test case, expected result, counterexample, or other verification 
artifact is modified during the main experiment.

The verification mechanism does not explicitly state whether the LLM recommendation 
is correct and does not provide an alternative LLM judgment. Instead, 
it provides task-specific evidence from which participants can assess 
whether the recommendation is consistent with the observed behaviour 
or the task requirements. The evidence therefore does not directly 
reveal the correctness label used by the experimenter.

Verification burden and verification cost are treated as distinct concepts. 
Verification cost refers to the intrinsic effort required to establish the 
correctness of a recommendation. Because this latent effort cannot be observed 
directly from participant behaviour---for example, a participant may stop 
before establishing correctness or may reach an incorrect conclusion---the 
study does not claim to measure intrinsic verification cost.

Instead, the experiment measures observable verification burden imposed by 
the system and experienced by the participant. The burden-related measures
are strictly separated by level:

\begin{itemize}
    \item \emph{verification availability}: whether verification evidence is made available to the participant (system-level);
    \item \emph{evidence examination}: whether the participant actually engages with the available evidence (participant-level);
    \item \emph{evidence examination time}: the time the participant spends interacting with the evidence panel;
    \item \emph{post-availability decision time}: the time from the moment evidence is made available to the final judgment, regardless of whether the evidence is examined (defined only on verification-available trials);
    \item \emph{total decision time}: the time from presentation of the LLM recommendation to the final judgment;
    \item \emph{post-evidence answer change}: whether the participant changes the initial answer after the evidence stage.
\end{itemize}

These terms are used consistently throughout the proposal: 
\emph{availability} refers to the system, \emph{examination} refers to the
participant, and the two time measures are distinct constructs. The term
\emph{verification time} is deliberately avoided because it conflates
system-level provision of evidence with participant-level use of evidence.

The experimental application records evidence-panel visibility events, 
activation and dismissal timestamps, scrolling and interaction events, 
and the timestamp of the final judgment. Evidence examination is 
operationalised using a rule specified before the main experiment: a trial 
is classified as \emph{evidence-examined} if the evidence panel remains 
active for at least $\delta$ seconds or if the participant performs a 
predefined interaction with the evidence panel. The value of $\delta$ is 
fixed during the pilot and frozen before the main experiment.

This instrumentation distinguishes system-level verification availability 
from participant-level evidence use. Evidence that is made available but 
not examined is still counted as verification availability, but is classified 
as non-examined at the participant level. Consequently, the adaptive condition 
can be evaluated not only by whether verification was made available less often 
than under static verification, but also by whether participants actually 
used the evidence when it was provided.

\subsection{Adaptive Verification Policy}
\label{subsec:policy}

The adaptive policy receives only the participant's initial confidence
 $c\in[0,100]$. Verification is made available when

\[
c < \tau,
\]

where $\tau$ is fixed before the main experiment. The threshold is either
specified a priori or selected using a separate pilot dataset. Pilot
participants are disjoint from main-experiment participants, and pilot data
are excluded from all main-experiment analyses; $\tau$ is not estimated,
optimised, or modified using data from the main experiment.

The policy does not receive the correctness of the LLM recommendation. Task
type and task difficulty are not policy inputs; they are retained as
experimental metadata and analysed only as exploratory factors (RQ3). The
policy therefore uses a single input---initial confidence---and nothing else.

The primary purpose of the policy is to test whether a simple
confidence-based gating rule can reduce verification availability relative to
an always-available (static) policy while maintaining decision quality. The
policy is deliberately minimal; its scientific value lies in the controlled
comparison against the always-available and no-verification baselines
(RQ1--RQ2), not in the rule itself.

\subsection{Policy Premise Check}
\label{subsec:premise}

The adaptive verification policy is based on the premise that initial
confidence provides information about the probability that the participant's
initial answer is correct. This premise is examined independently of the
policy decision. Specifically, initial-answer correctness is modelled as a
function of initial confidence using the mixed-effects structure specified in
Section~\ref{subsec:analysis}, with initial confidence scaled to the $[0,1]$ range and participant and task included as random effects.

The purpose of this analysis is to determine whether the assumed relationship
between confidence and initial-answer correctness is present in the collected
data. A positive association would provide empirical support for the
assumption that initial confidence contains information useful for
confidence-based gating: participants with lower initial confidence would, on
average, be more likely to hold an incorrect initial answer and could
therefore be more likely to benefit from verification. A weak or absent
association would indicate a boundary condition of the confidence-gating
approach, suggesting that initial confidence provides limited information for
selectively allocating verification.

This analysis does not determine or modify the verification threshold $\tau$.
The threshold is fixed before the main experiment according to the procedure
specified in the experimental design and remains unchanged throughout data
collection. The premise check is therefore an interpretive analysis rather
than a policy-tuning procedure.

The realised verification availability of the adaptive condition is reported
as an outcome of the fixed policy. In particular, the analysis reports the
proportion of trials for which the confidence threshold triggered verification
and examines how this availability relates to initial confidence. This allows
the study to assess whether the fixed policy actually reduces verification
availability relative to static verification without using the observed main
experiment data to optimise the policy.

\subsection{Participants, Recruitment, and Sample Size}
\label{subsec:participants}

\emph{Population.}

The target population is university students aged 18 or over who are enrolled
in or have completed at least one introductory programming course in any
programming language. The benchmark requires only basic Python reading,
control-flow, and input/output knowledge.

\emph{Inclusion criteria.}

Participants must:

\begin{enumerate}
\item be at least 18 years old;
\item have completed or be currently enrolled in an introductory programming
course, based on self-report;
\item demonstrate basic Python-reading competence using a short screening
quiz;
\end{enumerate}

\emph{Exclusion criteria.}

Participants are excluded if they:

\begin{enumerate}
\item fail the screening quiz according to the pre-registered criterion;
\item report prior exposure to the exact benchmark tasks; or
\item experience technical problems that prevent completion of the experiment.
\end{enumerate}

\emph{Recruitment.}

Participants are recruited through departmental mailing lists, course
announcements, and student groups. Compensation may consist of course credit
or a small voucher according to local institutional practice. Participation is
voluntary and requires informed consent. Participants may withdraw at any
time. Experimental logs are pseudonymised or anonymised as required by the
applicable institutional and data-protection procedures.

\emph{Sample size.}

The provisional target sample is $N=60$ participants, with a minimum feasible
sample of $N=36$. The final target sample size is determined by the
simulation-based power analysis described below and is fixed before
preregistration; the provisional values serve feasibility planning. 
With a provisional sample of \(N=60\), each participant contributes 18 trials,
yielding up to 1,080 trial-level records. These records are clustered 
within participants and tasks and are therefore not treated as independent 
observations.

The benchmark contains nine correct and nine incorrect LLM recommendations.
Task assignment is counterbalanced across participants as specified in
Section~\ref{subsec:answerspace} so that, within each condition,
approximately half of the trials contain correct recommendations and half
contain incorrect recommendations. Under the planned balanced assignment,
each condition therefore contributes approximately 180 observations with
correct recommendations and 180 observations with incorrect recommendations
at the provisional target sample size.

The six trials per condition are not treated as independent replications at
the participant level. Instead, the primary analysis uses trial-level
observations in mixed-effects models with participant and task included as
random effects. The sample size is therefore justified with respect to the
planned mixed-effects analysis rather than by treating the repeated trials
within a participant as independent observations.

Because each participant contributes only three trials with incorrect
recommendations and three trials with correct recommendations per condition,
participant-level estimates of over- and under-reliance are expected to be
unstable. These measures are therefore treated as trial-level binary
outcomes in the primary analysis, while aggregated rates are reported only
as descriptive summaries.

A simulation-based power analysis is conducted before preregistration using
the planned mixed-effects logistic models and variance components estimated
from the pilot data. The simulation varies plausible effect sizes, participant
variance, task variance, and the expected proportions of correct and incorrect
LLM recommendations. The final target sample size, together with the
confidence threshold $\tau$, the examination threshold $\delta$, and the
decision-time budget $B$, is fixed after the pilot and before main data
collection. Pilot participants are disjoint from main-experiment
participants, and pilot data are excluded from all main-experiment analyses.
Main-experiment data are not used to increase the sample size unless an
explicitly preregistered stopping or recruitment rule requires it.

A complete session is expected to take approximately 45--60 minutes,
including instructions, experimental trials, and the demographic
questionnaire.

\subsection{Experimental Design and Procedure}
\label{subsec:procedure}
Three conditions are compared:

\begin{enumerate}
\item \emph{No verification}: the participant receives the LLM recommendation
but no verification evidence is provided.
\item \emph{Static verification (always available)}: verification evidence is
always made available after the LLM recommendation.
\item \emph{Adaptive verification}: verification evidence is made available
only when the participant's initial confidence satisfies $c<\tau$.
\end{enumerate}

The verification mechanism is identical in all three conditions. Only the
display policy is manipulated.

Each participant completes all 18 benchmark tasks. The benchmark contains six
tasks of each type: \emph{identify bug}, \emph{check against specification},
and \emph{select correct implementation}. Within each task type, three tasks
contain correct LLM recommendations and three contain incorrect
recommendations.

Two tasks of each type are assigned to each experimental condition. Within
each condition, one task of each type has a correct LLM recommendation and one
has an incorrect LLM recommendation. Thus, each participant completes six
trials per condition, with three correct and three incorrect LLM
recommendations overall. Condition-to-task assignment is counterbalanced
across participants as specified in Section~\ref{subsec:answerspace}.

Each trial follows the same sequence:

\[
\begin{array}{c}
\text{Task}\\
\downarrow\\
\text{Initial answer}\\
\downarrow\\
\text{Confidence}\\
\downarrow\\
\text{LLM recommendation}\\
\downarrow\\
\text{Condition-specific verification}\\
\downarrow\\
\text{Final judgment}
\end{array}
\]

The participant first provides an answer independently. The participant then
reports confidence on a 0--100 scale. The fixed benchmark recommendation is
then displayed. Depending on the experimental condition, verification evidence
is either unavailable, always available, or made available according to the
confidence policy.

The participant subsequently records a final answer. The final answer may
confirm the initial answer or change it. The application does not reveal the
pre-determined correctness label of the recommendation. The correctness label
is stored only in the analysis dataset and is not available to the
participant-facing application or to the policy module during a trial.

For each trial, the system records the task identifier, task type, difficulty,
initial answer, initial confidence, fixed recommendation, condition,
adaptive-policy decision, recommendation correctness label for analysis only,
evidence availability, evidence-panel events, evidence examination status,
examination duration, post-availability decision time, final answer,
answer-change status, and relevant timestamps.

\subsection{Evaluation Metrics}
\label{subsec:metrics}

The evaluation metrics are organised into primary outcome measures, reliance
measures, and verification-related measures. The primary outcome measures are
final-answer correctness and reliance on the LLM recommendation. Reliance is
analysed at the trial level using the participant's acceptance of the LLM
recommendation, while over-reliance and under-reliance are derived from this
measure conditional on the objective correctness of the recommendation.

Final-answer accuracy is the proportion of trials on which the participant's
final structured answer is correct according to the benchmark reference
answer.

For each trial, LLM recommendation acceptance is defined as

$$
\text{AcceptLLM}_i =
I(\text{Final}_i=\text{LLM}_i),
$$

where equality is evaluated on the canonical structured response
representation defined in Section~\ref{subsec:answerspace}. Thus, acceptance
is determined exactly rather than by manual semantic comparison.

The primary statistical analysis models $\text{AcceptLLM}$ across all
trial-level observations using a mixed-effects logistic regression with
experimental condition, recommendation correctness, and their interaction as
fixed effects and participant and task as random intercepts. The
condition~$\times$~recommendation-correctness interaction is used to assess
whether the effect of verification differs between correct and incorrect LLM
recommendations. Condition contrasts among incorrect recommendations quantify
the effect on over-reliance, whereas condition contrasts among correct
recommendations quantify the effect on under-reliance.

For descriptive reporting, over-reliance and under-reliance are defined as
follows:

$$
\begin{array}{c}
\displaystyle
\text{OverReliance}
=
\frac{
N(\text{LLM incorrect}\land\text{Final}=\text{LLM})
}{
N(\text{LLM incorrect})
}
\\[1.5em]
\displaystyle
\text{UnderReliance}
=
\frac{
N(\text{LLM correct}\land\text{Final}\neq\text{LLM})
}{
N(\text{LLM correct})
}
\end{array}
$$

Over-reliance therefore represents the proportion of incorrect LLM
recommendations that are accepted, whereas under-reliance represents the
proportion of correct LLM recommendations that are rejected. These rates are
reported descriptively by experimental condition; inferential comparisons are
obtained from the primary mixed-effects model rather than from separate
subset models.

Appropriate reliance is the proportion of trials on which the participant
accepts a correct LLM recommendation or rejects an incorrect one:

$$
\begin{aligned}
\text{AppropriateReliance}
&=
\frac{
\begin{gathered}
N(\text{LLM correct}\land\text{Final}=\text{LLM})\\
{}+N(\text{LLM incorrect}\land\text{Final}\neq\text{LLM})
\end{gathered}
}{
N_{\text{trials}}
}.
\end{aligned}
$$

Decision-change rate is the proportion of trials on which the final answer
differs from the initial answer. Changes are classified along two independent
dimensions: (i) a correctness transition relative to the benchmark reference
answer (correct$\to$correct, correct$\to$incorrect, incorrect$\to$correct,
incorrect$\to$incorrect), which is used for outcome analysis; and (ii) the
relation to the LLM recommendation (toward or away from the recommendation),
which is reported descriptively. The two classifications are distinct. For
example, a change from a correct initial answer to a different incorrect
answer may move away from an incorrect recommendation while still producing
an incorrect final outcome.

\subsubsection{Verification Burden}
\label{subsubsec:verification-burden}

Verification burden is operationalised using directly observable measures.
The primary system-level measure is the \emph{verification availability
rate} (VAR):

$$
\text{VAR}
=
\frac{
N(\text{verification available})
}{
N_{\text{trials}}
}.
$$

For the No Verification condition, VAR is zero. For the Static Verification
condition, VAR is one apart from pre-specified exclusions. For the Adaptive
Verification condition, VAR depends on the participant's confidence and the
fixed threshold $\tau$.

Participant-level burden is characterised by:

\begin{itemize}
\item evidence-examination rate;
\item evidence-examination time;
\item post-availability decision time;
\item total decision time;
\item proportion of displayed evidence that was not examined;
\item proportion of trials followed by a final-answer change.
\end{itemize}

These measures are reported separately rather than combined into an aggregate
burden score.

\subsubsection{Verification Success}
\label{subsubsec:verification-success}

Verification success is evaluated relative to a pre-specified decision-time
budget $B$. The budget is an analysis-time censoring threshold on the
post-availability decision time rather than an interface timeout: participants
are not forced to stop at $B$. The value of $B$ is calibrated using the pilot,
for example near the 90th percentile of observed post-availability decision
times, and is frozen before the main experiment.

Let $V_i$ denote the participant's implicit verdict about the LLM
recommendation on trial $i$:

$$
V_i =
\begin{cases}
\text{accept}, & \text{if Final}_i=\text{LLM}_i,\\
\text{reject}, & \text{otherwise}.
\end{cases}
$$

Let $C_i$ indicate whether the LLM recommendation is objectively correct. A
verdict is considered correct when the participant accepts a correct
recommendation or rejects an incorrect recommendation.

Let $E_i$ indicate that verification evidence was available and examined on
trial $i$, as defined in Section~\ref{subsec:burden}, and let
$D_i^{\mathrm{post}}$ denote the post-availability decision time. The
verification success indicator is

$$
S_{B,i}
=
I\left(
V_i\text{ is correct}
\land
D_i^{\mathrm{post}}\leq B
\right),
$$

evaluated only on trials with $E_i=1$. The verification success rate is

$$
\text{VSR}
=
\frac{
N(S_B=1)
}{
N(E=1)
}.
$$

VSR is therefore a conditional measure characterising whether participants
who engaged with the verification evidence reached a correct verdict within
the decision-time budget. Trials in which evidence was available but not
examined are excluded from the primary VSR because the evidence was not
actually engaged with. The same statistic computed over all
verification-available trials is reported as a descriptive supplement.

VSR does not establish that verification caused the correct decision, because
a participant may already have held the correct judgment before examining the
evidence.

For incorrect recommendations, the verification failure rate is

$$
\text{VFR}_{\mathrm{incorrect}}
=
\frac{
N(\text{LLM incorrect}\land E=1\land S_B=0)
}{
N(\text{LLM incorrect}\land E=1)
}.
$$

This measures the proportion of incorrect recommendations that were not
rejected within the decision-time budget despite the evidence being available
and examined.

Verification success and verification burden are reported separately. No
aggregate quality--burden score is constructed.

\subsection{Statistical Analysis}
\label{subsec:analysis}

The primary analyses are conducted at the trial level using mixed-effects
regression models. Repeated observations from the same participant and the
same benchmark task are not treated as independent. Participant and task are
therefore included as random effects to account for participant-level and
task-level heterogeneity. Because task--condition assignments are
counterbalanced across participants (Section~\ref{subsec:answerspace}), task
characteristics are not systematically associated with a particular
experimental condition.

For binary outcomes, mixed-effects logistic regression models are used. The
general model specification is

$$
\operatorname{logit}\!\left(P(Y_{ij}=1)\right)
=
\beta_0
+
\boldsymbol{\beta}^{\mathsf{T}}\mathbf{X}_{ij}
+
u_i
+
v_j,
$$

where $Y_{ij}$ is a binary outcome for participant $i$ on task $j$,
$\mathbf{X}_{ij}$ represents the fixed-effect predictors, $u_i$ is a
participant-level random intercept, and $v_j$ is a task-level random
intercept.

The primary predictor is experimental condition, with No Verification as the
reference condition. Planned contrasts compare Adaptive Verification with No
Verification and Static Verification with No Verification. A further
comparison between Adaptive and Static Verification evaluates whether the
adaptive policy produces comparable outcomes while making verification
available on fewer trials.

Recommendation correctness is included as a planned factor for reliance
analyses because the effect of verification is expected to differ between
correct and incorrect LLM recommendations. The interaction between
experimental condition and recommendation correctness is therefore a primary
term in the reliance analysis. Task type and task difficulty may be included
as exploratory moderators. Neither task type nor task difficulty is used as
an input to the adaptive verification policy, and analyses involving these
variables are reported as exploratory.

The primary outcome is final-answer correctness, modelled at the trial level
using a mixed-effects logistic regression with experimental condition as the
primary predictor and participant and task as random intercepts. Planned
condition contrasts assess differences in final-answer correctness between
the experimental conditions.

Reliance on the LLM recommendation is analysed using a single mixed-effects
logistic regression on all trial-level observations. For each trial,
acceptance of the LLM recommendation is defined as

$$
\text{AcceptLLM}_{ij}
=
I(\text{Final}_{ij}=\text{LLM}_{j}),
$$

where equality is evaluated using the canonical structured response
representation defined in Section~\ref{subsec:answerspace}. The primary
reliance model is

$
\operatorname{logit}\!\left(
P(\text{AcceptLLM}_{ij}=1)
\right)
=
\beta_0
+
\beta_1\text{Condition}_{ij}
+
\beta_2\text{RecCorrect}_{j}
+
\beta_3
\left(
\text{Condition}_{ij}
\times
\text{RecCorrect}_{j}
\right)
+
u_i
+
v_j.
$

The condition-by-recommendation-correctness interaction tests whether the
effect of verification on acceptance differs between correct and incorrect
LLM recommendations. For incorrect recommendations, acceptance corresponds
to over-reliance; therefore, the planned condition contrasts within this
level quantify differences in over-reliance. For correct recommendations,
rejection corresponds to under-reliance; therefore, the planned condition
contrasts within this level quantify differences in under-reliance. This
single model is preferred to separate regression models fitted to correct and
incorrect recommendation subsets because it uses all trial-level observations
and estimates the condition-by-correctness interaction directly.

Aggregate over-reliance, under-reliance, and appropriate-reliance rates are
reported descriptively by experimental condition. Over-reliance is the
proportion of incorrect recommendations that are accepted, whereas
under-reliance is the proportion of correct recommendations that are
rejected. Appropriate reliance is the proportion of trials on which the
participant accepts a correct recommendation or rejects an incorrect
recommendation.

The trial-level indicator of appropriate reliance may additionally be
analysed as a secondary mixed-effects logistic outcome. Answer-change rate is
analysed separately as a binary outcome indicating whether the final answer
differs from the initial answer. Where appropriate, answer changes are also
classified by their correctness transition relative to the benchmark
reference answer.

The primary analysis of verification availability compares the proportion of
trials in which verification is made available under the Static and Adaptive
conditions. Because verification availability is determined directly by the
fixed adaptive policy, this analysis describes the realised behaviour of the
policy rather than estimating or optimising its threshold. Verification
availability in the No Verification condition is structurally zero and is
therefore not included in the Static-versus-Adaptive comparison.

The relationship between initial confidence and initial-answer correctness is
examined separately as the policy premise check described in
Section~\ref{subsec:premise}. Initial-answer correctness is modelled using a
mixed-effects logistic regression:

$$
\operatorname{logit}\!\left(
P(Y^{\mathrm{init}}_{ij}=1)
\right)
=
\beta_0
+
\beta_1 c_{ij}
+
u_i
+
v_j,
$$

where $Y^{\mathrm{init}}*{ij}$ indicates whether the initial answer of
participant $i$ on task $j$ is correct and $c*{ij}$ is the reported initial
confidence scaled to the $[0,1]$ range. The coefficients $u_i$ and $v_j$
represent participant- and task-level random intercepts, respectively.
Confidence is scaled before analysis; no further transformation is applied.
This analysis evaluates whether initial confidence contains information about
initial-answer correctness without using the result to modify the adaptive
threshold $\tau$.

The realised verification availability of the adaptive condition is also
reported as a descriptive outcome. In particular, the proportion of adaptive
trials for which the fixed confidence threshold triggers verification is
reported together with its distribution across initial-confidence levels.
These quantities are not used to optimise $\tau$ using main-experiment data.

Time-based outcomes, including post-availability decision time and total
decision time, are analysed separately because their distributions may be
non-normal and right-skewed. Where appropriate, log-transformed times or
generalized mixed-effects models with an appropriate response distribution
are used. Trials with technically invalid timing measurements are excluded
according to a predefined data-cleaning rule. Analyses of verification
examination time and post-availability decision time are restricted to trials
on which verification evidence was made available.

The verification success rate is treated as a secondary, conditional outcome
because it is defined only for trials on which verification evidence was
available and examined. Its interpretation and calculation are specified in
Section~\ref{subsubsec:verification-success}. The analysis does not treat
verification success as an independent causal effect of verification.

All primary hypotheses, contrasts, model specifications, exclusion rules, and
coding of outcome variables are defined before analysis of the main
experimental data. Pilot participants and pilot observations are excluded
from all main-experiment analyses. The significance level and procedure for
reporting confidence intervals and effect sizes are specified in the
preregistration. Exploratory analyses are reported separately from the
primary confirmatory analyses.


\subsection{Expected Software Artefacts}

The thesis will deliver:

\begin{itemize}

\item a frozen benchmark containing 18 introductory Python tasks;

\item structured task metadata including type, difficulty, reference answer,
fixed LLM recommendation, and recommendation correctness;

\item the archived Qwen3-8B 8-bit model specification, exact model artifact
identifier/hash, inference configuration, prompts, candidate generations, and
validation records;

\item the experimental web application;

\item the execution-based verification module;

\item the configurable confidence-based policy module;

\item structured experimental logging including evidence availability,
evidence examination, examination time, post-availability decision time, and
total decision time;

\item the frozen benchmark-generation and validation scripts;

\item the statistical analysis scripts implementing the pre-registered
mixed-effects models; and

\item documentation sufficient to reproduce benchmark generation,
experimental execution, and statistical analysis.

\end{itemize}


%=====================================================================
\section{Detailed Work Plan and Schedule}
\label{sec:workplan}
%=====================================================================


The work is planned from October 2026 to mid-May 2027, with the main
objective of submitting a research paper by 15 May 2027; the thesis is then
extended with additional documentation, analysis, and discussion.


\begin{longtable}{p{8.0cm}p{2.4cm}p{5.3cm}}
\caption{Work plan and schedule}
\label{tab:workplan} \\
\toprule
\textbf{Phase and activities} & \textbf{Period} &
\textbf{Deliverable / milestone} \\
\midrule
\endfirsthead


\toprule
\textbf{Phase and activities} & \textbf{Period} &
\textbf{Deliverable / milestone} \\
\midrule
\endhead


\bottomrule
\endfoot


1. Literature review and research design:
systematic study of appropriate reliance, verification cost, human--AI
interaction, confidence-aware support, and LLM assistance in programming;
completion of the bibliography &
October 2026 &
Annotated bibliography; refined research gap and research questions \\


2. Finalisation of the experimental design:
operational definitions, hypotheses, metrics, burden--cost distinction,
confidence-based policy, statistical analysis plan, and power-simulation
setup &
1--15 November 2026 &
Draft protocol and analysis plan \\


3. Benchmark construction and LLM recommendations:
selection of 18 introductory Python tasks, categorisation, difficulty
assessment, reference-answer validation, offline generation and validation of
recommendations (9 correct / 9 incorrect), and task-specific verification
evidence &
15--30 November 2026 &
Validated benchmark with frozen recommendations (M2) \\


4. Prototype implementation:
participant interface, task manager, fixed recommendation module, verification
module, adaptive policy, counterbalancing procedure, and experiment logger
including evidence-examination logging &
December 2026 &
Functional experimental prototype (M3) \\


5. Ethics, pilot, and protocol freeze:
ethics approval; pilot recruitment and pilot testing of task clarity,
interface usability, logging quality, confidence elicitation, and calibration
of the confidence threshold $\tau$, the examination threshold $\delta$, and
the decision-time budget $B$; simulation-based power analysis;
pre-registration &
January 2027 &
Frozen configuration ($\tau$, $\delta$, $B$, $N$) and registered protocol (M4) \\


6. Main experiment:
recruitment until the preregistered target sample is reached (minimum
analysable $N=36$), counterbalanced assignment, data collection, monitoring
of exclusions, quality control &
Late January--February 2027 &
Complete and anonymised trial dataset (M5) \\


7. Data analysis:
data cleaning, descriptive statistics, mixed-effects models for accuracy and
reliance, burden and examination analysis, premise check, trade-off
comparison of static and adaptive verification &
February--March 2027 &
Reproducible analysis scripts and preliminary results (M6) \\


8. Paper preparation:
writing the methods, results, discussion, limitations, and conclusion;
tables and figures; internal revision and proofreading &
March--April 2027 &
Complete paper draft and revised manuscript (M7) \\


9. Paper submission:
final revision, formatting, reference checking, supplementary materials,
submission &
By 15 May 2027 &
Submitted research paper (M8) \\


10. Thesis continuation and integration:
expansion of the paper into the thesis, documentation of the software
artefacts, extended analysis, final thesis version &
After 15 May 2027 &
Extended thesis manuscript and project documentation \\


\end{longtable}

%=====================================================================
\section{References}
\printbibliography[heading=none]


\end{document}